\section{실패 자산화와 회귀 테스트}
탐색에서 찾은 실패는 한 번의 영상으로 끝나지 않고 재검증 가능한 테스트 자산이 되어야 한다.
모든 실행은 장면 파라미터, 시뮬레이션 모델, 작업 목표, 모델·소프트웨어 버전, 영상, 상태·접촉 기록,
로봇 기능 실행 결과를 해시값과 함께 저장한다. 회귀 실행기는 이 기록을 읽어 동일
조건의 사례를 구성하고, 현재 버전으로 정해진 횟수 안에서 재실행한 뒤 원본과 비교한다.

\begin{enumerate}
\item 저장된 PASS/FAIL 사례를 읽고 동일 조건·장면의 사례를 묶는다.
\item 사례별 반복 수와 호출 상한을 정하고 입력·자산 호환성을 검사한다.
\item 현재 버전 조건을 기록한 뒤 재실행한다.
\item 원본과 신규 결과를 표~\ref{tab:regstate}의 상태로 분류하고 영상·근거가 연결된 보고서를 생성한다.
\end{enumerate}

\begin{table}[!t]
\centering\footnotesize\setlength{\tabcolsep}{3pt}
\caption{회귀 비교 상태}\label{tab:regstate}
\begin{tabularx}{\linewidth}{@{}p{26mm}Y@{}}
\toprule
상태 & 의미 \\
\midrule
유지 (성공 / 실패) & 이전과 같은 작업 성공 여부 \\
퇴행 & 이전에 성공한 사례가 새 버전에서 실패 \\
개선 & 이전에 실패한 사례가 새 버전에서 성공 \\
혼재 & 이전 또는 새 반복에서 성공과 실패가 섞임 \\
대기 / 판정 불가 & 미완료 또는 제외로 비교 불가 \\
\bottomrule
\end{tabularx}
\end{table}

\begin{table}[!t]
\centering\footnotesize\setlength{\tabcolsep}{3pt}
\caption{로봇 LLM 교체 후 휴머노이드 사례의 회귀 재실행 결과 (본 실험 장면)}\label{tab:regrun}
\begin{tabularx}{\linewidth}{@{}lYcc@{}}
\toprule
원래 사례 & 재실행 조건 & 결과 & 분류 \\
\midrule
시작 장면 PASS & Qwen3.8-27B, 3회 & PASS 2·제외 1 & 판정 불가 \\
(Luna 7/7) & 추가 1회 & PASS 1 & 유지(성공) \\
목표 점유 FAIL & Qwen3.8-27B, 3회 & FAIL 3 & 유지(실패) \\
목표 점유 FAIL & GPT-6 Luna, 2회 & FAIL 2 & 유지(실패) \\
\bottomrule
\end{tabularx}
\end{table}

표~\ref{tab:regrun}은 로봇 LLM을 GPT-6 Luna에서 Qwen3.8-27B로 바꾸어 저장된 본 실험 사례를 재실행한 결과다(장면·호출
상한·평가 규칙 동일). 시작 장면은 유효 실행이 모두 성공했고, API 호출 한도 오류(HTTP 429)로 끝난 1회는 실패가 아닌 판정 불가로
분리되었다. 목표 점유 실패는 두 모델 모두에서 유지되어, 이 사례들에서는 모델 교체로 결과가 바뀌지 않았다. 다만 Qwen은 시작 장면을 끝내는 데 로봇 판단을 중앙값 10회(5·10·13회) 사용하여, 같은 장면의 Luna 7회(중앙값 5회, 4--9회)의
약 2배였다. 결과 분류는 유지지만 판단 상한에 더 가까워진 것으로, 회귀 테스트는 PASS/FAIL 외에 판단 횟수·소요 시간 같은
자원 사용량 변화도 함께 추적해야 함을 보여 준다. 로봇팔 개발 단계
점검에서도 23개 환경을 3회씩 재실행해 원래 FAIL 10건이 모두 재현되었다.
본 과제에서 발견한 장애물 높이 경계(12/14\,cm)와 목표 점유 장면은 정책이나 현장 배치가 바뀔 때 우선 재실행할
회귀 테스트셋의 핵심 사례가 된다.
